\documentclass{article}

% TMLR format
\usepackage[accepted]{tmlr}

\newcommand{\openreview}{\url{https://openreview.net/forum?id=XXXX}}
\def\month{06}
\def\year{2026}

\usepackage{amsmath,amssymb,amsthm}
\usepackage{booktabs}
\usepackage{multirow}
\usepackage{hyperref}
\usepackage{xcolor}
\usepackage{graphicx}
\usepackage{enumitem}
\usepackage{microtype}

\newcommand{\TODO}[1]{{\color{orange}[TODO: #1]}}
\newcommand{\NUM}[1]{\textbf{#1}}
\newcommand{\PENDING}[1]{{\color{blue}[PENDING: #1]}}

\title{\textbf{Phonological Assembly Circuits in GPT-2:\\
How Transformers Compute Sound Structure from Orthography}}

\author{
  Elliot Tower \\
  \texttt{elliot@elliottower.com}
}

\begin{document}

\maketitle

\begin{abstract}
No published work has performed circuit-level causal analysis on phonological
operations in transformer language models. We identify and validate circuits for
six distinct phonological operations in GPT-2 Small, ranging from stored lexical
lookup (rhyming hypocorism: Richard $\to$ Dick) to cross-boundary phoneme fusion
(oronym: ``Bill Ding'' $\to$ ``building''). Each operation requires a
linguistically distinct computation --- irregular lookup, regular truncation,
grapheme-to-phoneme mapping, juncture reanalysis, homophone recognition, and
reverse folk etymology --- and we find that each activates a \emph{distinct}
circuit architecture with \PENDING{N--M} heads per circuit.

Using weight-space circuit discovery \citep{tower2026weights}, EAP, and
activation patching, we show that: (1)~the universal backbone tier (L0H8, L0H9,
L0H11) participates in all phonological circuits, consistent with the
``operating system'' hypothesis; (2)~the oronym circuit --- which must fuse
phonemes across a word boundary --- uses \PENDING{specific mid-layer heads
in layers 3--8} that have not been identified in any prior circuit;
(3)~compound phonological puns (``Richard Tater'' $\to$ ``Dictator'') activate
the expected composition of component circuits, and ablating either component
breaks the compound.

The oronym circuit is the central novel finding: it demonstrates that GPT-2
implements cross-boundary phonological assembly from purely orthographic input,
a capability previously demonstrated only behaviorally
\citep{liao2026tokenization,suvarna2024phonologybench} but never localized to
specific attention heads.
\end{abstract}

%% =========================================================
\section{Introduction}
%% =========================================================

The mechanistic interpretability literature has characterized circuits for
syntactic operations (IOI, SVA, reflexive anaphora) and simple retrieval tasks
(induction, greater-than). But the space of \emph{phonological} circuits ---
computations that operate on sound structure rather than syntactic structure ---
is entirely unexplored at the circuit level.

This gap matters because phonological computation in transformers is
surprising: BPE tokenization encodes orthography, not phonology, yet language
models demonstrably encode phonological structure in their hidden states
\citep{liao2026tokenization}. Linear probes on GPT-2's middle layers achieve
63--79\% accuracy on rhyme detection, and larger models perform above chance
on grapheme-to-phoneme conversion \citep{suvarna2024phonologybench}. The
question we address is: \emph{which specific attention heads compute this
phonological knowledge, and do different phonological operations use different
circuits?}

\paragraph{The phonetic name pun as a natural task.}
We use phonetic name puns as our task family because they decompose into
well-specified linguistic operations with clear minimal-pair evaluations.
``Richard Tater $\to$ Dictator'' requires two operations: (Op~1) rhyming
hypocorism (Richard $\to$ Dick) and (Op~4) oronym/juncture reanalysis
(Dick Tater $\to$ Dictator). By testing each operation independently, we
isolate the circuit for each computation.

\paragraph{Contributions.}
\begin{enumerate}[leftmargin=*,itemsep=2pt]
  \item Six linguistically specified circuit tasks with minimal-pair datasets
  (\PENDING{total N} pairs across 6 operations).
  \item Circuit discovery for each operation using EAP, activation patching,
  and weight-space features.
  \item The first identification of cross-boundary phoneme fusion heads in
  any transformer model (the oronym circuit).
  \item Cross-reference with the GPT-2 acronym circuit
  \citep{garcia2024acronym}: Op~3 (initialism) \PENDING{shares / does not
  share} heads with their letter-mover circuit.
  \item Compositional validation: compound puns activate the expected
  combination of component circuits.
\end{enumerate}

%% =========================================================
\section{Operations and Datasets}
%% =========================================================

\subsection{Linguistic Classification}

Each operation is defined by the linguistic computation it requires. We use
standard terminology from phonology and morphology.

\begin{table}[h!]
\centering
\caption{The six phonological operations tested, with linguistic classification,
example pairs, and dataset sizes.}
\label{tab:operations}
\small
\begin{tabular}{llllr}
\toprule
\textbf{Op} & \textbf{Linguistic term} & \textbf{Example} & \textbf{Computation} & \textbf{N pairs} \\
\midrule
1 & Rhyming hypocorism & Richard $\to$ Dick & Irregular stored lookup & 32 \\
2 & Clipping (apocope) & Timothy $\to$ Tim & Regular first-syllable truncation & 90 \\
3 & Initialism + letter name & J $\to$ Jay & Grapheme $\to$ phoneme name & 30 \\
4 & Oronym (juncture reanalysis) & Bill Ding $\to$ building & Cross-boundary phoneme fusion & 99 \\
5 & Homophone recognition & Neil $\to$ kneel & Phoneme-identity detection & 120 \\
7 & Folk etymology & microwave $\to$ Michael Wave & Reverse decomposition & 75 \\
\bottomrule
\end{tabular}
\end{table}

\subsection{Dataset Construction}

Each pair has the format \texttt{\{good, bad, target\_token, foil\_token\}},
following the BLiMP minimal-pair methodology \citep{warstadt2020blimp}. The
evaluation metric is logit difference: the model should assign higher
probability to the target than the foil at the final position.

\paragraph{Tokenization audit.}
Following \citet{liao2026tokenization}, we audit every pair for BPE
tokenization: target and foil must each be a single BPE token.
\PENDING{N/total} pairs pass this filter. Multi-token targets are excluded
from circuit discovery but reported as a tokenization bottleneck finding.

\paragraph{Procedural expandability.}
Ops 2, 3, and 5 can be expanded procedurally using CMU Pronouncing Dictionary
(Op~2: syllable-based clipping, Op~3: all 26 letters, Op~5: $\sim$3000
homophone pairs). Ops 1, 4, and 7 are controlled vocabularies --- acceptable
for psycholinguistic-style studies as long as the full inventory is reported
\citep{warstadt2020blimp}.

%% =========================================================
\section{Behavioral Validation}
%% =========================================================

Before circuit discovery, we test whether GPT-2 Small can perform each
operation at all. Gate: $>$70\% accuracy on logit-diff evaluation.

\begin{table}[h!]
\centering
\caption{Behavioral validation results. Gate threshold: 70\% accuracy.}
\label{tab:behavioral}
\small
\begin{tabular}{lrrrl}
\toprule
\textbf{Operation} & \textbf{Accuracy} & \textbf{Mean LD} & \textbf{N} & \textbf{Gate} \\
\midrule
Op 1: Rhyming hypocorism   & \PENDING{xx\%} & \PENDING{$\pm$x.xx} & 32  & \PENDING{PASS/FAIL} \\
Op 2: Clipping             & \PENDING{xx\%} & \PENDING{$\pm$x.xx} & 90  & \PENDING{PASS/FAIL} \\
Op 3: Letter name          & \PENDING{xx\%} & \PENDING{$\pm$x.xx} & 30  & \PENDING{PASS/FAIL} \\
Op 4: Oronym               & \PENDING{xx\%} & \PENDING{$\pm$x.xx} & 99  & \PENDING{PASS/FAIL} \\
Op 5: Homophone            & \PENDING{xx\%} & \PENDING{$\pm$x.xx} & 120 & \PENDING{PASS/FAIL} \\
Op 7: Folk etymology       & \PENDING{xx\%} & \PENDING{$\pm$x.xx} & 75  & \PENDING{PASS/FAIL} \\
\bottomrule
\end{tabular}
\end{table}

\paragraph{Expected results.}
Ops 1 and 2 (name $\to$ nickname) should pass easily --- these are well-known
associations in GPT-2's training data. Op 5 (homophones) should pass based on
PhonologyBench results. Op 4 (oronym) is uncertain --- this requires
cross-boundary phoneme fusion, which is bottlenecked by tokenization. Op 7
(folk etymology) is expected to fail on GPT-2 Small since it requires
generating creative decompositions not in training data.

\paragraph{Cross-model sweep.}
For operations that fail the gate on GPT-2 Small, we test GPT-2 Medium and
Large. \PENDING{Report which ops pass at which scale.}

%% =========================================================
\section{Circuit Discovery}
%% =========================================================

For each operation passing the behavioral gate, we run three complementary
discovery methods:

\subsection{EAP (Edge Attribution Patching)}

Gradient-based head importance via retained gradients on \texttt{hook\_z}
activations. Ranks all 144 heads by attribution score.

\subsection{Activation Patching}

Single-head mean ablation: replace each head's activations with corpus mean
and measure logit-diff degradation.

\subsection{Weight-Space Features}

Extract 173 weight features (108 spectral + 65 behavioral proxy) per head.
Bootstrap stability classification with greedy formula search identifies the
weight-predicted circuit \citep{tower2026weights}.

\subsection{Results}

\PENDING{Per-operation tables with: (a) EAP top-15 ranking, (b) activation
patching top-15, (c) weight-predicted circuit, (d) P/R/F1 of weight prediction
vs activation ground truth.}

\subsubsection{Op 1: Rhyming Hypocorism Circuit}

\PENDING{Expected: lookup circuit. Small number of heads with high OV cosine
similarity between name input direction and nickname output direction. Likely
concentrated in layers 3--8 (mid-layer, where Liao \& Shi find phonological
info peaks). Should resemble IOI name-mover heads but for associative content
rather than positional copying.}

\begin{table}[h!]
\centering
\caption{Op 1 circuit: top heads by activation patching.}
\label{tab:op1_circuit}
\small
\begin{tabular}{lrr}
\toprule
\textbf{Head} & \textbf{Patching effect} & \textbf{Weight-predicted?} \\
\midrule
\PENDING{L?H?} & \PENDING{x.xxx} & \PENDING{Y/N} \\
\PENDING{L?H?} & \PENDING{x.xxx} & \PENDING{Y/N} \\
\PENDING{...}   &                  &                \\
\bottomrule
\end{tabular}
\end{table}

\subsubsection{Op 2: Clipping Circuit}

\PENDING{Expected: positional/syllable-boundary detector circuit. Related to
acronym circuit's positional heads but operating on within-token phonological
structure.}

\subsubsection{Op 3: Initialism + Letter Name Circuit}

\PENDING{This is the micro-version of the acronym task. Key test: do the
Garc\'ia-Carrasco et al. 8 letter-mover heads appear in this circuit?
If yes, convergent validation with independent method.
If no, Op 3 requires additional phoneme-lookup heads beyond positional
extraction, which would be a new finding.}

\subsubsection{Op 4: Oronym Circuit (Core Novel Finding)}

\PENDING{This is the most novel circuit. The model must fuse phonemes across
a word boundary (``Bill'' + ``Ding'' $\to$ ``building''). No prior work has
identified cross-boundary phoneme fusion heads.

Expected circuit structure:
\begin{itemize}
  \item Backbone (L0): universal infrastructure
  \item Mid-layer fusion heads (L3--L8): the novel component --- heads that
  attend across the word boundary and compose phonemic representations
  \item Readout (L10--L11): standard output infrastructure
\end{itemize}

Key question: does this circuit overlap with the clipping circuit (Op 2)?
If the fusion heads are distinct from clipping heads, that's evidence for
specialized phonological circuitry.}

\subsubsection{Op 5: Homophone Circuit}

\PENDING{Expected: phoneme-identity detection circuit in mid-layers.
Liao \& Shi's probing results (63--79\% rhyme detection accuracy) suggest
phonological representations exist in layers 3--8. The homophone circuit
should read from these representations.}

\subsubsection{Op 7: Folk Etymology Circuit}

\PENDING{Expected to fail behavioral gate on GPT-2 Small. If it passes on
Medium/Large, the circuit would be the reverse of Op 4 --- decomposing a word
into name-like constituents rather than composing names into a word.}

%% =========================================================
\section{Cross-Reference Analysis}
%% =========================================================

\subsection{Shared Infrastructure}

\PENDING{Does the backbone tier (L0H8, L0H9, L0H11) appear in all phonological
circuits, consistent with the ``operating system'' hypothesis from Paper A?
Expected: yes, based on transfer results across RTI/SVA/IOI/gendered pronoun.}

\subsection{Overlap with Acronym Circuit}

The Garc\'ia-Carrasco et al.\ acronym circuit consists of 8 letter-mover heads.
Op 3 (initialism) is the closest match to their task.

\begin{table}[h!]
\centering
\caption{Overlap between phonological circuits and known circuits.}
\label{tab:overlap}
\small
\begin{tabular}{lccccc}
\toprule
& \textbf{RTI} & \textbf{IOI} & \textbf{SVA} & \textbf{Acronym} & \textbf{Induction} \\
\midrule
Op 1 (hypocorism) & \PENDING{N} & \PENDING{N} & \PENDING{N} & \PENDING{N} & \PENDING{N} \\
Op 2 (clipping)   & \PENDING{N} & \PENDING{N} & \PENDING{N} & \PENDING{N} & \PENDING{N} \\
Op 3 (letter name)& \PENDING{N} & \PENDING{N} & \PENDING{N} & \PENDING{N} & \PENDING{N} \\
Op 4 (oronym)     & \PENDING{N} & \PENDING{N} & \PENDING{N} & \PENDING{N} & \PENDING{N} \\
Op 5 (homophone)  & \PENDING{N} & \PENDING{N} & \PENDING{N} & \PENDING{N} & \PENDING{N} \\
\bottomrule
\end{tabular}
\end{table}

\subsection{The Phonological Circuit Overlap Matrix}

\PENDING{Heat map showing pairwise overlap between all 6 operation circuits.
Expected: Ops 1 and 2 (both name-shortening) share more heads than Ops 4 and 5
(different mechanisms). If Op 4 and Op 5 share mid-layer heads, that suggests
a shared ``phonological representation'' layer.}

%% =========================================================
\section{Compositional Validation}
%% =========================================================

\PENDING{Test compound puns that chain operations:
\begin{itemize}
  \item ``Richard Tater'' $\to$ ``Dictator'' = Op 1 (Richard $\to$ Dick) + Op 4 (Dick Tater $\to$ Dictator)
  \item ``Elizabeth Aardvark'' $\to$ ``Lizard'' = Op 2 (Elizabeth $\to$ Liz) + Op 4 (Liz Ard $\to$ Lizard)
  \item ``William Dingo'' $\to$ ``Building'' = Op 1 (William $\to$ Bill) + Op 4 (Bill Ding $\to$ Building)
\end{itemize}

For each compound: (a) test behavioral accuracy, (b) ablate Op 1 circuit only
--- compound should break, Op 4 alone should survive, (c) ablate Op 4 circuit
only --- compound should break, Op 1 alone should survive. This is the
compositionality test.

NOTE: This section requires running compound ops after individual op circuits
are found. Not included in initial pod run.}

%% =========================================================
\section{Tokenization as Bottleneck}
%% =========================================================

\PENDING{Following Liao \& Shi (2026):
\begin{itemize}
  \item Report STAD scores for all input words
  \item Test slash-delimited condition (``B/i/l/l D/i/n/g'') for Op 4
  \item If slash-delimiter significantly improves Op 4 accuracy, that's
  evidence that the oronym circuit \emph{exists} but is bottlenecked by
  tokenization --- the model has the structural capacity (weight-space
  circuit) but can't activate it with standard tokenization
  \item This would be a publishable finding connecting the weight-capacity
  vs.\ activation-flow distinction from Paper A to a phonological setting
\end{itemize}}

%% =========================================================
\section{Discussion}
%% =========================================================

\paragraph{Novelty.}
No prior work has done circuit-level causal analysis on phonological
operations. The closest is the acronym circuit \citep{garcia2024acronym},
which is positional (letter extraction) rather than phonological. Our work
connects the behavioral phonological capabilities demonstrated by
\citet{suvarna2024phonologybench} and \citet{liao2026tokenization} to specific
attention heads for the first time.

\paragraph{The oronym circuit as the core finding.}
Cross-boundary phoneme fusion --- composing two input tokens (``Bill'' +
``Ding'') into a phonologically coherent output (``building'') --- is the
most novel circuit type. If GPT-2 can do this, the model has learned to
represent phonological continuity across token boundaries despite training
on purely orthographic input.

\paragraph{Implications for the operating system hypothesis.}
If the backbone tier participates in all phonological circuits, it extends
the ``operating system'' finding from Paper A/D to a new domain. The backbone
provides general-purpose attention infrastructure; phonological computation is
built on top of it, not in parallel.

\paragraph{Limitations.}
\begin{itemize}[leftmargin=*,itemsep=2pt]
  \item Controlled vocabulary for Ops 1, 4, 7 (not procedurally generated).
  \item GPT-2 Small may fail behavioral gate on hardest operations.
  \item English-only; phonological operations differ across languages.
  \item BPE tokenization is a known confounder --- results reflect the
  interaction of phonological knowledge with tokenization artifacts.
\end{itemize}

\paragraph{Venue target.}
COLM, Findings of ACL, or MI workshop. The novel operation types,
compositional structure, and first-ever phonological circuit identification
position this as a meaningful extension of the MI circuits literature.

\bibliographystyle{tmlr}
\begin{thebibliography}{99}

\bibitem[Tower(2026)]{tower2026weights}
Tower, E. (2026). Reading circuits from weights. \textit{Manuscript}.

\bibitem[Garc\'ia-Carrasco et~al.(2024)]{garcia2024acronym}
Garc\'ia-Carrasco, J. et~al. (2024). How does GPT-2 predict acronyms?
  \textit{AISTATS 2024}. arXiv:2405.04156.

\bibitem[Suvarna et~al.(2024)]{suvarna2024phonologybench}
Suvarna, A., Khandelwal, A., and Peng, N. (2024). PhonologyBench: Evaluating
  phonological skills of large language models. arXiv:2404.02456.

\bibitem[Liao and Shi(2026)]{liao2026tokenization}
Liao, C. and Shi, P. (2026). How tokenization limits phonological knowledge
  representation in language models. \textit{ACL 2026}. arXiv:2604.17105.

\bibitem[Warstadt et~al.(2020)]{warstadt2020blimp}
Warstadt, A. et~al. (2020). BLiMP: The benchmark of linguistic minimal pairs
  for English. \textit{TACL}.

\bibitem[Wang et~al.(2023)]{wang2023ioi}
Wang, K. et~al. (2023). Interpretability in the wild. \textit{ICLR 2023}.

\bibitem[Olsson et~al.(2022)]{olsson2022icl}
Olsson, C. et~al. (2022). In-context learning and induction heads.
  \textit{Transformer Circuits Thread}.

\bibitem[Hanna et~al.(2023)]{hanna2023gt}
Hanna, M. et~al. (2023). How does GPT-2 compute greater-than?
  \textit{NeurIPS 2023}.

\end{thebibliography}

\end{document}
